\chapter{Runtime contracts and request lifecycles}
A model step finishes just as its client cancels. May the engine publish the
result? May it reuse the step's state buffer for another tenant before the
device acknowledges completion? What if the output stream still contains two
committed tokens? These are semantic questions before they are performance questions.

Chapter 29 froze evidence and comparison contracts. This chapter constructs the
serving shell: immutable request identity, admission, exclusive state leases,
bounded output, commit fences, terminal arbitration and cleanup. The reference
does not execute DOGMA or HermonDNA mathematics. It specifies the lifecycle that
later native execution must preserve. A correct shell is not a trained genomic
model, a network service or evidence of industrial throughput.

\section{Bind a request before allocating state}
Represent a request by
\[
 R=(r,u,m,d,N,Q_I,Q_B,B_{\max}).
\]
$r$ is request identity, $u$ tenant/security scope, $m$ a pinned model-artifact
identity, and $d$ a deadline in a declared monotonic clock domain.
$N$ bounds total committed output items; $Q_I$ bounds queued items;
$Q_B$ bounds queued UTF-8 bytes; $B_{\max}$ bounds one item.
These integer ceilings are independent. An empty queue does not mean the request
has unused generation budget.
\par\medskip\noindent\begin{minipage}{\linewidth}
\Listing{Request}
\end{minipage}\par\medskip
The identity format is lowercase 64-character SHA-256. Chapter 28 checked pinned
bytes for its linear artifact fixture; a production engine must extend that
binding to the actual model, configuration, vocabulary, precision and semantic
options. A correctly shaped string does not authenticate an artifact. The reference
assumes a trusted binding supplied by its caller.
\Plate{01-envelope}{The envelope binds independent fields; arrows do not pretend
that tenant identity causes model identity. Policy and availability checks precede
state publication. Failed admission leaves no expensive state lease.}{evod30:envelope}

A hot deployment may choose a new identity for \emph{new} requests.
An existing request does not follow a mutable current-model pointer.
Migrating its native state needs an explicit, verified conversion protocol.
Matching tensor shape alone does not establish state compatibility.

\section{Admission is a decision, not a scheduler}
The oracle requires an explicit policy result and model-availability result.
Before its deadline, an allowed request becomes admitted; a policy or availability
failure becomes rejected. Deadline expiry terminates without allocation.
Scheduling subsequently chooses execution order and placement among admitted
work. Chapter 34 owns scheduling and batching.

Admission must account for resource reservation if concurrent requests compete.
A check-then-allocate sequence without synchronized reservation can overbook a pool.
Our registry's exclusivity contract prevents state-ID aliasing, not global GPU
capacity exhaustion. There is no memory allocator, quota database, authentication
service or scheduling fairness proof here. Unknown available capacity is not
permission to treat it as infinite.

\section{Terminal outcome and cleanup are separate coordinates}
Use phases received, admitted, active and terminal. Streaming is a property of
an active or terminal request's queue, not a phase that erases ownership.
Keep a terminal reason and a separate cleaned flag:
\[
 S=(\phi,\rho,c,\ell,v,p,\mathcal Q,\mathcal A).
\]
$\phi$ is phase, $\rho$ terminal reason, $c$ cleanup state, $\ell$ the lease,
$v$ the committed revision, $p$ the pending ticket, $\mathcal Q$ the queue and
$\mathcal A$ accounting counters.
\Plate{02-lifecycle}{Lifecycle phase, immutable terminal reason and resource
cleanup are separate. A completed request may still have queued output.
A cancelled request may still have physical work awaiting acknowledgement.}
{evod30:lifecycle}
Replacing completed by cleaned loses the outcome unless it is stored elsewhere.
Similarly, reporting cancellation as complete because some bytes streamed hides
a truncated operation. The reference retains exactly one reason from completed,
cancelled, deadline, failed or rejected. Cleanup does not overwrite it.

\section{One serialized terminal winner}
All operations below are stipulated to execute serially.
A terminal transition is the \Term{linearization point} for outcome selection:
before it, eligible work may commit; after it, no model result may create a new
committed output or revision.
\Listing{Runtime.terminate}
Before the deadline, the first eligible terminal proposal wins.
At or after the deadline, a newly observed terminal transition records deadline
instead of pretending to complete on time. Already-terminal requests retain their
earlier winner. Equality at $d$ means expired, not one last allowed commit.

The six permutations of completion, cancellation and failure before expiry are
tested against an independent first-winner rule. Four proposed reasons at exactly
$d$ all select deadline. Completion requires active state and no pending work;
cancellation may terminate a received or admitted request too.
This is an explicit arbitration policy, not a rule forced on every inference API.

For a concurrent implementation, serialize the \emph{whole} check/update protocol
or prove a suitable atomic transaction, including lease, deadline, revision,
queue and counters. An atomic terminal flag alone cannot protect an unguarded
output append. A test with twelve threads under one external lock checks terminal
serialization only; the reference itself is not advertised as thread-safe.

\section{Ownership survives physical placement changes}
The registry maps each live state ID to
\[
 \ell=(s,g,(u,r,m)),
\]
where $g$ is a monotonically increasing allocation generation for that ID.
A request owns its logical state; a worker is a residency label.
The registry rejects even a second claim by the same owner while the ID is live.
A shared prefix or fork is not created accidentally by repeated claim.
\Plate{03-ownership}{The same physical state ID is released and later receives a
new generation. A completion carrying the old generation must not write into the
new tenant's allocation. Worker placement is separate from logical ownership.}
{evod30:ownership}

Allocation claims the registry entry before publishing active local state.
An ownership conflict leaves the request admitted and unallocated.
A foreign release fails. After release and reuse, the old lease also fails even
if the new allocation has the same request/model names.
This counters the \Term{ABA reuse} problem: the string S may look unchanged while
the resource it denotes has changed.

Generations persist only in this in-memory oracle. Across a crash or worker restart,
a production design needs globally unique allocation epochs or a durable generation
scheme and fencing at the executor. These tokens are correctness identities, not
cryptographic capabilities. The registry assumes trusted code; it cannot prevent an
attacker mutating process memory.

\section{Backpressure before model work}
To avoid silently running a step whose result has nowhere to go, reserve a worst-case
output slot before launch:
\[
 |\mathcal Q|<Q_I,\qquad
 \operatorname{bytes}(\mathcal Q)+B_{\max}\le Q_B.
\]
At most one ticket is pending for this request.
This simple policy is conservative: a two-byte item in a five-byte queue leaves
three bytes free, but another item with declared maximum four bytes cannot launch.
\Listing{Runtime.launch}
The blocked path logs backpressure but creates no ticket and increments no launched
work counter. A producer cannot repeatedly advance hidden state or consume sampling
RNG while calling those attempts blocked output. A native engine must make state
and RNG commit part of the same operation or preserve a pending snapshot.
\Plate{04-stream}{Item and byte ceilings both constrain launch. Total committed
output is a third limit. A slow consumer holds only bounded buffered bytes; the
backpressure path does not perform another model step.}{evod30:stream}

UTF-8 character count is not byte count. The item \emph{é} encodes to two bytes;
a supplementary-plane emoji commonly needs four. The reference stores immutable
encoded bytes and validates each produced item's maximum. It does not claim to
bound the original Python string's transient allocation, native state, event log,
allocator overhead or all process memory. The trace itself is unbounded here;
a production system needs retention, sampling or bounded/durable export policies.

\section{Fence a late result, then commit atomically}
A launch ticket binds the exact lease, current revision and unique launch number:
\[
 p=(\ell,v,j_{\rm launch}).
\]
A result must match the pending ticket and a still-live lease. Duplicate,
foreign and stale tickets cannot settle a different step.
\par\medskip\noindent\begin{minipage}{\linewidth}
\Listing{Runtime.commit}
\end{minipage}\par\medskip
After deadline/control checks, a successful result advances revision and output
count together and appends one ordered byte payload. A result after a terminal
winner settles the physical-work ticket but increments discarded-work accounting;
it does not advance state or queue.
Invalid type, empty output, malformed Unicode or oversized output terminates as
failed with no new revision. The oracle does not model a fault during event append
or crash halfway through the Python mutation sequence; production atomicity must
cover such failures separately.

\Plate{05-races}{A launched step can finish physically after cancellation.
Its acknowledgement settles the ticket, but the output/state commit is fenced.
The lease cannot be released while physical work may still access its buffer.}
{evod30:races}
Cancellation is therefore not instantaneous device preemption.
If a kernel still holds a pointer, logically cancelling the request does not make
that memory safe to reassign. A timeout waiting for acknowledgement is an unresolved
resource, not a license to reuse the buffer. Device synchronization, abort guarantees
or a quarantined allocation are required.

\section{Deadlines use one explicit clock domain}
Times in the reference are exact nonnegative integer milliseconds, passed by tests.
Time never decreases. All admission, allocation, launch, commit, consumption and
control paths observe the deadline where relevant. A step launched before $d$
but returning at $d$ cannot commit a timely result.

Production intervals should use a monotonic clock, not civil time that may jump.
Monotonic epochs on different hosts are not automatically comparable. Transmit
a deadline using a documented synchronization or remaining-budget protocol,
include transit/error bounds, and check the receiving clock domain.
A local service cannot guarantee a remote client's receipt time merely by meeting
its own commit deadline.

HTTP/3 distinguishes rejecting an unprocessed request from cancelling one whose
application processing has begun \cite{evod30:http3}. That protocol distinction
motivates careful retry semantics; it does not define this model's state machine.

\section{Draining is not success, and cleanup is not delivery}
Our chosen queue policy keeps previously committed bytes readable after every
terminal reason until cleanup. No new model outputs can commit. A client must
receive the terminal reason separately to distinguish partial output from success.
A different API may discard on cancellation, but it must declare and count that loss.

For the three-output fixture, the first two commits fill the queue; a launch is
blocked. Consuming item one allows item three to commit and automatically reach
the total output ceiling. The completed request still has two buffered items.
\begin{center}\input{\Generated/stream.tex}\end{center}
Sequence numbers increase one through three regardless of consume timing.
A local dequeue is not proof of remote receipt, application processing or exactly-once
delivery. Transport reconnect and acknowledgement need their own protocol.

\par\medskip\noindent\begin{minipage}{\linewidth}
\Listing{Runtime.cleanup}
\end{minipage}\par\medskip
Cleanup requires terminal phase and no pending physical work. It also requires
an empty queue or an explicit discard decision. The registry releases the lease
before local handles are cleared. A repeated successful cleanup returns a no-op.
If release fails, local handles are retained so the failure remains diagnosable.
This is logical idempotence in one live process, not a crash-resilient finalizer.

\section{Reconcile every counter with an invariant}
Let $L,S,P$ be launched, settled and pending work, with $P\in\{0,1\}$.
Let $G,Q,D,X$ be committed outputs, queued, locally dequeued and discarded outputs.
Then
\[
 L=S+P,\qquad G=Q+D+X,\qquad v=G,\qquad G\le N.
\]
Discarded \emph{work} is distinct from discarded \emph{output}.
A cancelled in-flight step can be settled and discarded without ever creating
a committed output; a committed but undelivered item can be discarded during cleanup.
\Plate{06-trace}{The trace records launches, commits, local dequeues, terminal
selection and cleanup as distinct events. Separate work and output conservation
identities expose a missing acknowledgement or an unreported queue loss.}
{evod30:trace}

The normal fixture finishes with $L=S=G=D=v=3$, $P=Q=X=0$.
The cancelled pending fixture finishes with $L=S=1$, discarded work one, $G=v=0$.
Counters measure declared logical events, not FLOPs, device time, billable tokens
or a general resource cost. Billing and scheduler reservations require additional
accounting contracts.

\section{Independent interleavings before load tests}
The test alphabet is launch, commit, consume, cancel, complete and cleanup.
Enumerate all $6^5=7776$ length-five sequences from an allocated request.
For each prefix, compare accepted/no-op/rejected operations and final state with a
separate small specification that does not call the runtime's methods.
Check queue bounds, one terminal event, lease release and both conservation
identities after every operation.

This is bounded exhaustive verification for a stipulated serialized oracle.
It is not exploration of every distributed execution, weak-memory ordering or
device race. Extend the alphabet with worker crashes, loss/reorder, scheduler
reservation changes, model-step failure and durable recovery before claiming
production correctness. Chapter 26's restricted affine-scan checks do not supply
complete DOGMA/HermonDNA engine parity; those model-specific checks remain distinct.

The trace contains immutable event records with monotone sequence numbers and
times, but its list is in memory. It omits output payloads and hidden activations.
It is neither an authenticated audit log nor a replay-complete durable journal.
A production design must reconcile transactional publication and trace persistence,
including crash windows, privacy, retention and recovery.

\section{Retries and execution authority are additional contracts}
A generation request is not automatically idempotent because it uses HTTP.
RFC 9110 distinguishes method-level idempotence and unsafe automatic retries
\cite{evod30:http}. A client retry might reconnect to one retained operation or
create a new one. Durable tenant-scoped idempotency keys must bind normalized
request content and retention policy; a cleaned request ID alone is insufficient.
If the same key accompanies different content, reject rather than silently reuse.

An inference shell also does not authorize actions. For an agent integration,
model output proposes; a policy/kernel authorizes; a deterministic service executes;
a durable trace records. Tool execution needs separate capabilities, revocation,
side-effect idempotence and recovery. Neither a token nor a runtime trace creates
execution authority or proves intelligence.

\section{Twenty worked exercises}
\begin{enumerate}
\item Why pin $m$? \textbf{Answer:} avoid mid-request model drift; identity is not authentication.
\item Is $Q_I=N$ required? \textbf{Answer:} no; queue occupancy and total output are different ceilings.
\item Why check availability before allocation? \textbf{Answer:} rejected work should obtain no native-state lease.
\item Does admission schedule the request? \textbf{Answer:} no; it decides eligibility.
\item Can completed imply cleaned? \textbf{Answer:} no; buffered output and physical resources may remain.
\item Where is the terminal linearization point? \textbf{Answer:} the serialized transition that fixes phase and reason.
\item Completion proposed exactly at $d$? \textbf{Answer:} deadline wins if the request was not already terminal.
\item Cancellation after successful completion? \textbf{Answer:} no reason rewrite or resurrection.
\item What does a generation number fix? \textbf{Answer:} stale leases after state-ID reuse, including same-name ABA.
\item Why isn't it a security capability? \textbf{Answer:} trusted callers and process integrity are assumed.
\item Queue bytes two, ceiling five, item maximum four: launch? \textbf{Answer:} no; worst-case reservation would require six bytes.
\item How many bytes is é in UTF-8? \textbf{Answer:} two; one character does not mean one byte.
\item A pending ticket after cancellation: release its buffer? \textbf{Answer:} no; acknowledge physical quiescence first.
\item A late result's state revision? \textbf{Answer:} unchanged; settled/discarded work increases instead.
\item Why retain committed bytes after cancellation? \textbf{Answer:} it is this API's explicit drain policy, not universal semantics.
\item Is local dequeue remote delivery? \textbf{Answer:} no; transport/application acknowledgement is separate.
\item Derive output conservation. \textbf{Answer:} each committed item is queued, dequeued or explicitly discarded, so $G=Q+D+X$.
\item What does $6^5$ verification establish? \textbf{Answer:} bounded serialized-sequence agreement, not all concurrent schedules.
\item Is the event list durable? \textbf{Answer:} no; storage/recovery/privacy contracts remain unimplemented.
\item Design a serving release gate. \textbf{Rubric:} native-step parity, pinned identity,
tenant ownership, synchronized commit/terminal protocol, clock domain, queue and
resource bounds, physical quiescence, fault injection, durable trace/retry protocol,
isolation, measured workload and bounded performance claims.
\end{enumerate}

\section{Bridge to native state}
Chapter 31 constructs and advances model-native state inside this shell.
A ticket's committed revision must correspond to an actual legal prefix and a
verified state transition; it cannot merely increment a counter. Full DOGMA carry
and HermonDNA KV semantics remain separate, and Chapter 21's suffix-aware offline
representation must not enter causal left-to-right generation.
